% Planteamiento del problema, justificación, objetivos y alcance. El problema
% sigue la indicación del tutor: la necesidad de un simulador para redes
% descentralizadas, formulada como una sola pregunta.

Este capítulo formula el problema que da origen al trabajo, justifica la
implementación de una herramienta propia y fija los objetivos y el alcance.

\section{Planteamiento del problema}
\label{sec:problema}

Quien quiera evaluar un protocolo para una red descentralizada dispone de tres
tipos de herramienta, y cada uno deja por fuera algo que aquí resulta esencial
(Capítulo~\ref{cap:estado-arte}). Los simuladores de propósito general, como
ns-3~\cite{Riley2010ns3} u OMNeT++~\cite{Varga2008Omnetpp}, evalúan una copia del
protocolo escrita para la herramienta, de modo que lo que se comprueba no es el
programa que se instalaría en un equipo. Los emuladores, como
CORE~\cite{Ahrenholz2008Core} o Mininet-WiFi~\cite{Fontes2015MininetWifi},
ejecutan el programa real, pero cada nodo consume recursos del computador y no
incluyen un modelo de la pérdida de señal entre los pisos de un edificio. Las
herramientas para redes entre pares, como PeerSim~\cite{Montresor2009Peersim},
trabajan en la capa de aplicación y no representan el medio de radiofrecuencia
ni el movimiento físico de los nodos.

A estas diferencias se suma un riesgo menos visible. Un simulador conoce, por
construcción, todo el sistema: las posiciones de los nodos, los enlaces y el
camino más corto entre cualquier par. Si el experimento usa esa información para
decidir algo que en la red real decidiría un nodo, ya no está evaluando una red
descentralizada, en la que cada nodo solo dispone de información
local~\cite{Corson1999RFC2501}. El simulador del que parte este
trabajo~\cite{JofsanchezciSimulador}, concebido para otro propósito, usaba esa
información global en tres lugares: el movimiento de los Gateway consultaba las
posiciones de todos los nodos, los mensajes se entregaban buscando el camino en
la red real, con garantía de llegada si existía alguno, y no había paquetes de
datos que pusieran a prueba las tablas de rutas. Además, la calidad de
transmisión (TQ), la medida con la que el protocolo elige sus
rutas~\cite{Seither2011BatmanAdv}, valía siempre 1.0, porque nunca se
registraban los mensajes perdidos.

Por otra parte, en una red de este tipo las preguntas de interés no se
responden con una sola cifra de conectividad: qué cree cada nodo y cuánto se
aleja eso de la realidad, cuánto tarda la red en notar un fallo y adaptarse, qué
nodos son críticos y si la red puede reorganizarse sin que nadie la coordine.
La literatura sobre la confianza en los estudios de simulación pide, además, que
cada resultado se pueda repetir y que se declare qué parte del sistema es un
modelo~\cite{Pawlikowski2002Credibility,Kurkowski2005Incredibles}.

El problema se formula entonces con la siguiente pregunta:

\begin{quote}
\textit{¿Cómo implementar un simulador para redes descentralizadas que ejecute
el código de enrutamiento de los nodos, represente el medio de radiofrecuencia
de una edificación, tome sus decisiones solo con información local y permita
medir la cobertura y la adaptabilidad de la red?}
\end{quote}

En este documento, la \textit{cobertura} es la fracción de usuarios que la malla
alcanza, y la \textit{adaptabilidad} es la capacidad de la malla para
reorganizarse por sí misma cuando cambian las condiciones, por ejemplo cuando un
nodo se cae o se mueve: rehacer sus rutas, volver a unirse y recuperar a los
usuarios que quedaron sin servicio. Ambas se miden con las métricas de la
Sección~\ref{sec:metricas}.

\section{Justificación}
\label{sec:justificacion}

La justificación es de tres tipos. En lo \textbf{científico}, la literatura ha
señalado problemas de confianza en los estudios de simulación de
redes~\cite{Kurkowski2005Incredibles,Pawlikowski2002Credibility} y que los
resultados dependen de la herramienta
empleada~\cite{Cavin2002Accuracy}. Una herramienta en la que cada corrida se
define con un archivo y una semilla, y en la que se declara qué parte es
programa real y qué parte es modelo, atiende directamente esos problemas.

En lo \textbf{práctico}, las mallas inalámbricas se usan para dar conectividad
donde no hay infraestructura o esta ha fallado~\cite{Akyildiz2005Mesh,Reina2015Disaster}.
Poder explorar con bajo costo su comportamiento antes de desplegarlas (dónde
poner los equipos, qué pasa si se cae uno, cuántos hacen falta) es útil para
quien las diseña, y los bancos de prueba físicos son costosos y difíciles de
repetir~\cite{Kiess2007RealWorld}.

En lo \textbf{académico}, la interfaz permite ver lo que un nodo cree, su tabla
de rutas y los mensajes con los que se propaga esa información, lo que convierte
a la herramienta en un recurso para enseñar cómo funcionan estas redes. Además,
las reglas locales de movimiento relacionan este trabajo con la inteligencia
artificial: los Gateway son agentes que, sin coordinación central y con
información parcial, producen un comportamiento colectivo, como en los métodos de
despliegue autónomo~\cite{Howard2002PotentialFields,Poduri2004ConstrainedCoverage}.
Esto deja un entorno de pruebas para aprender esas reglas
(Sección~\ref{sec:recomendaciones}).

\section{Objetivos}
\label{sec:objetivos}

\subsection{Objetivo general}

Implementar un simulador para redes descentralizadas que ejecute la lógica de
enrutamiento B.A.T.M.A.N. de los nodos sobre un modelo del medio de
radiofrecuencia y permita estudiar, de forma repetible y solo con decisiones
locales, la cobertura y la adaptabilidad de una Red de Expansión de Cobertura
ante fallos de nodos.

\subsection{Objetivos específicos}
\label{subsec:objetivos-especificos}

\begin{enumerate}
\item \textbf{Requerimientos.} Definir los requerimientos del simulador a partir
de la revisión de las herramientas de simulación y emulación existentes.
\item \textbf{Diseño.} Diseñar la arquitectura del simulador, el modelo del medio
de radiofrecuencia y los mecanismos con los que los nodos deciden solo con
información local.
\item \textbf{Implementación.} Implementar el simulador con escenarios
configurables, envío de paquetes de datos, movimiento de los Gateway, caídas
programables de nodos y una interfaz para observarlo y manejarlo.
\item \textbf{Pruebas.} Probar el simulador con simulaciones de una Red de
Expansión de Cobertura que midan su cobertura y su adaptabilidad ante distintos
tipos de fallo.
\end{enumerate}

El Capítulo~\ref{cap:desarrollo} sigue este mismo orden.

\section{Alcance y limitaciones}
\label{sec:alcance}

\textbf{Se incluye:} un simulador en Python que avanza en pasos de 0.5~s, con
uso por terminal, ventana propia e interfaz web local; escenarios en un edificio
por pisos y en campo abierto; caídas y recuperaciones programables, al azar o
dirigidas; tráfico de datos configurable con paquetes de tipo UDP, ICMP y TCP;
tres modos de movimiento de los Gateway y tres de los usuarios; y el análisis de
resultados por lotes de corridas.

\textbf{No se incluye:} el despliegue físico de la malla; la comparación contra
el módulo batman-adv del núcleo de Linux~\cite{OpenMesh2024Batmanadv} o contra
otro simulador; las colisiones, la interferencia y el detalle de la capa de
acceso al medio; la medición del retardo, porque cada salto se resuelve dentro de
un paso; ni la seguridad de la red.

Dos precisiones condicionan la interpretación del trabajo. Primero, el código
que se ejecuta es la implementación en Python del nodo de red de este proyecto
(paquete \texttt{mesh/}), no el módulo batman-adv, y difiere de este en puntos
que se documentan en la Sección~\ref{sec:fidelidad}. Segundo, el modelo del medio
es una función simple de la distancia y de los pisos, de modo que hereda
supuestos que se saben optimistas frente a un canal
real~\cite{Kotz2004Assumptions}. Por ello, los resultados deben leerse como el
comportamiento de esta implementación bajo este medio, y no como una predicción
numérica de un despliegue físico.
